% 由 analysis/ingest_multiclass.py 自动生成 —— 勿手改
% 由 analysis/ingest_multiclass.py 自动生成 —— 勿手改
% 来源: experiments/real_system/results/seeded_multiclass.json (run_multiclass.py)
% 真值 = 靶机正确模型 intended_perm；观测 = 实际请求；与真实系统同一套 alarm 公式。
\newcommand{\McInstances}{9}
\newcommand{\McReadHits}{3}
\newcommand{\McChannelHits}{5}
\newcommand{\McOracleReadHits}{7}
\newcommand{\McOracleChannelHits}{9}
\newcommand{\McVthreeReadHits}{3}
\newcommand{\McVthreeChannelHits}{5}
\newcommand{\McFpBookkeeping}{8}
\newcommand{\McAuthorizedCells}{23}
\newcommand{\McFpOracle}{0}
\newcommand{\McClasses}{4}


\begin{table*}[t]
\centering
\caption{Per-defect recall of the studied chain on the controlled target, by defect class. Each instance is a cell whose true (correct-model) permission is \emph{none} and on which the defect nevertheless lets the access succeed. \emph{read gate} is the reachability-gated relation as studied in this paper; \emph{channel gate} is the same relation moved to the channel the defect actually breaks; the last two columns are the oracle-consuming adjudicator on the read channel and on the channel-matched probe. Every instance here is one the gate was designed for except C3 and C4, which it cannot see: in both, the object is visible to the subject, so the antecedent of the relation is false and no alarm is raised---not a tuning problem but the gate reading its own premise. The oracle-based adjudicator reaches all \McInstances{} instances once the probe is channel-matched. False alarms on the \McAuthorizedCells{} authorized cells: \McFpBookkeeping{} for the relation's bookkeeping variant, \McFpOracle{} for the oracle-based adjudicator, whose zero is conditional on this target's oracle being sound (Section~VI-C reports what happens when it is not).}
\label{tab:multiclass}
\small
\begin{tabular*}{\textwidth}{@{\extracolsep{\fill}}lrrrrr}
\toprule
Class & inst. & read gate & channel gate & oracle (read) & oracle (channel) \\
\midrule
C1: read enforcement & 3 & 3 & 3 & 3 & 3 \\
C2: modify enforcement & 2 & 0 & 2 & 0 & 2 \\
C3: inheritance & 2 & 0 & 0 & 2 & 2 \\
C4: privilege escalation & 2 & 0 & 0 & 2 & 2 \\
\addlinespace[2pt]
\emph{total} & 9 & 3 & 5 & 7 & 9 \\
\bottomrule
\end{tabular*}
\end{table*}
